\section{Experimentación y Resultados}

\subsection{Dataset Final}

\subsection{Implementación del/los Modelo/s}

\subsubsection{Extractor/core}

A continuación se detalla el \textit{pipeline} completo que inicia con un archivo de correo electrónico (de extensión
\texttt{.eml}) y que finaliza con la extracción de los datos relevantes del mismo.

Para el caso de correos electrónicos correspondientes a vuelos se busca obtener el código de la reserva, el nombre del
pasajero y el detalle de uno o más segmentos de vuelo.
Cada segmento se compone de un código de vuelo, el aeropuerto de origen, el aeropuerto de destino y la fecha y hora de
inicio de vuelo.

\paragraph{Archivo \texttt{.eml}, HTML y texto a simple vista}

El primer paso consiste en obtener el contenido del correo electrónico y quedarse con solo lo relevante.
Para ello se utiliza la librería simple-java-email (\url{https://www.simplejavamail.org}) que convierte el contenido de
un \texttt{.eml} a formato HTML. Para ello se utiliza el método \texttt{org.simplejavamail.api.email.Email\#getHTMLText}
y se consigue una instancia de \texttt{String} con todo el contenido HTML\@.

A continuación se utiliza la librería Jsoup (\url{https://jsoup.org/}) que consiste en un \textit{parser} html que
ofrece el método \texttt{org.jsoup.nodes.Element\#text()} para aplanar el HTML del mail y quedarse únicamente con el
texto visible, tal como si se estuviera visualizando el mail desde una aplicación de correo electrónico.
Esta acción aplica también para todas las secciones posibles según el estándar MIME Multipart
(\url{https://en.wikipedia.org/wiki/MIME}).
Un correo electrónico cuenta con varias de estas secciones, algunas contienen adjuntos, otras contienen información de
reenvío del mail, y algunas incluso codifican el texto en Base64 (\url{https://en.wikipedia.org/wiki/Base64}) de forma
que la biblioteca procede a decodificarlo.

Es importante destacar que el \textit{pipeline} trabaja únicamente con el texto visible de los mails, ignorando mucho
más contenido que podría ser relevante para identificar directamente áreas críticas del mail.
Atributos presentes en los \textit{tags} html, imágenes, adjuntos como archivos PDF y demás contenidos son ignorados
por el \textit{pipeline}.
Incluso el asunto del mail es ignorado.

Además, se aplica un preprocesamiento extra para eliminar \textit{footers} conocidos. Por ejemplo aquellos nodos HTML
identificados con cadenas como \texttt{.footer} son ignorados.
Durante la implementación este paso fue relevante para ignorar por ejemplo información de contacto de empresas de
aerolíneas que pueden confundir al algoritmo NLP (un caso es el nombre del presidente de la compañía, el cual se podría
confundir con el nombre del pasajero).
Dado que está comprobado que no existe información útil para el usuario ubicada en esta sección entonces es seguro
eliminarla.

Otro preprocesamiento adicional que se aplica es la búsqueda mediante expresiones regulares de cadenas que puedan
corresponder con códigos de vuelo.
Un código de vuelo IATA (\url{https://en.wikipedia.org/wiki/Airline_codes}) consiste en dos letras, seguido de un número
de entre 2 y 4 cifras.
Por ejemplo \texttt{AC8633} corresponde al vuelo 8633 de la aerolínea Air Canada.
En caso de que el mail contenga cadenas como \texttt{AC-8633} o \texttt{AC 8633} este preprocesamiento elimina el guion
o el espacio para que siempre se obtenga la cadena \texttt{AC8633}, facilitando así la tarea de la próxima etapa del
\textit{pipeline} para identificar códigos de vuelo usando una única convención.

\begin{lstlisting}[
    style=javacodestyle,
    caption={Normalización de códigos de vuelo.},
    label={lst:normalize-flight-numbers}
]
private static String normalizeFlightNumbers(String text) {
  Pattern FLIGHT_NUMBER_WITH_HYPHEN_SEPARATOR =
    Pattern.compile("\\b([A-Z][A-Z0-9]|[0-9][A-Z]|[A-Z][0-9])\\s*[-–]\\s*(\\d{2,4})\\b");
  Pattern FLIGHT_NUMBER_WITH_SPACE_SEPARATOR =
    Pattern.compile("\\b([A-Z][A-Z0-9]|[0-9][A-Z]|[A-Z][0-9])\\s+(\\d{2,4})\\b");
\end{lstlisting}

\paragraph{Stanford CoreNLP}

Una vez que se cuenta con una instancia de \textit{String} con únicamente el contenido a simple vista de un correo
electrónico se envía este texto al \textit{pipeline} de Stanford CoreNLP.
El mismo cuenta con múltiples etapas que detallamos a continuación:

\subparagraph{Tokenization}

(\url{https://stanfordnlp.github.io/CoreNLP/tokenize.html})

La tokenización consiste en convertir texto en tokens.
La forma más elemental consiste en dividir cada vez que aparece un espacio.
Por ejemplo para el texto ``Thanks for flying.'' la tokenización genera cuatro tokens: ``Thanks'', ``for'', ``flying'' y
``.''.
CoreNLP realiza la separación del texto en \textit{tokens} usando una colección de reglas ya previstas.
En el archivo de configuración se indica la propiedad \texttt{tokenize.language=en} para que utilice las reglas del
idioma inglés.
Desde este momento la implementación estará limitada exclusivamente a correos electrónicos de reservas redactados en
inglés.

\subparagraph{Sentence Splitting}

(\url{https://stanfordnlp.github.io/CoreNLP/ssplit.html})

La separación en oraciones consiste en convertir un texto en oraciones.
Por ejemplo ``Thanks for flying. Flight safe.''.
será separado en las oraciones ``Thanks for flying.'' y ``Flight safe.''.
Aquí también CoreNLP realiza la separación en oraciones usando una colección de reglas ya previstas.

\subparagraph{Parts of Speech}

(\url{https://stanfordnlp.github.io/CoreNLP/pos.html})

Partes del discurso asignan \textit{tags} o etiquetas a los \textit{tokens} para marcarlos según su categoría
gramatical.
Por ejemplo para el texto ``Marie was born in Paris.'' se etiqueta a ``Marie'' como \texttt{NNP}
(nombre propio singular), ``was'' como \texttt{VDB} (verbo pasado), ``born'' como \texttt{VBN} (verbo participio pasado),
``in'' como \texttt{IN} (preposición), ``París'' nuevamente como \texttt{NNP} y el punto final ``.'' como signo de
puntuación.

\subparagraph{Lemmatization}

(\url{https://stanfordnlp.github.io/CoreNLP/lemma.html})

Aquí se asocia una palabra con su lema o forma de diccionario.
Por ejemplo para la palabra ``was'' se mapea por la palabra ``be''.
Este paso lleva todo el contenido a un solo tiempo verbal para simplificar la tarea de las próximas etapas.

\subparagraph{Named Entity Recognition}

(\url{https://stanfordnlp.github.io/CoreNLP/ner.html})

En esta etapa, la más compleja de todo el pipeline de NLP, se reconocen entidades en el texto como personas o nombres
de empresas.
La librería utiliza uno o más modelos de secuencia de \textit{machine learning} para etiquetar las entidades.
Además permite la ejecución de reglas específicas para etiquetar, por ejemplo, fechas, tiempos y valores numéricos.

Para el idioma inglés la biblioteca reconoce por defecto 24 tipos distintos de entidades:
\begin{itemize}
    \item Nombres reconocidos: \texttt{PERSON}, \texttt{LOCATION}, \texttt{ORGANIZATION}, \texttt{MISC}
    \item Numéricos: \texttt{MONEY}, \texttt{NUMBER}, \texttt{ORDINAL}, \texttt{PERCENT}
    \item Temporales: \texttt{DATE}, \texttt{TIME}, \texttt{DURATION}, \texttt{SET}
    \item Otros: \texttt{EMAIL}, \texttt{URL}, \texttt{CITY}, \texttt{STATEORPROVINCE}, \texttt{COUNTRY},
    \texttt{NATIONALITY}, \texttt{RELIGION}, \texttt{TITLE}, \texttt{IDEOLOGY}, \texttt{CRIMINALCHARGE},
    \texttt{CAUSEOFDEATH}, \texttt{HANDLE}
\end{itemize}

A simple vista se puede apreciar que no existe una entidad aeropuerto o aerolínea por lo que es necesario crear reglas
adicionales para facilitar el reconocimiento de nuevos tipos de entidades.

\subsubparagraph{Aeropuertos}

Un aeropuerto de pasajeros se identifica por un código de tres letras denominado \textit{IATA Code}
(\url{https://en.wikipedia.org/wiki/IATA_airport_code}).
A excepción de aeropuertos muy pequeños, cualquier aeropuerto cuenta con un identificador único de tres letras emitido
por IATA.

Para esta implementación se utilizó el \textit{dataset} público de Our Airports Data
(\url{https://ourairports.com/data/}) que ofrece la descarga de un archivo CSV con más de 84.000 registros que
representan a todos los aeropuertos del mundo.
Se procede a filtrar el mismo para considerar únicamente los aeropuertos que tengan un código IATA (descartando
aeropuertos pequeños).
Por ejemplo el aeropuerto ``Montreal / Pierre Elliott Trudeau International Airport'' se identifica con el código IATA
``YUL''.

\begin{lstlisting}[
    style=basecodestyle,
    caption={Registro de ejemplo del \textit{dataset} Our Airports Data.},
    label={lst:ourairports-example}
]
id,ident,type,name,latitude_deg,longitude_deg,elevation_ft,continent,iso_country,iso_region,municipality,scheduled_service,icao_code,iata_code,gps_code,local_code,home_link,wikipedia_link,keywords
1928,CYUL,large_airport,Montreal / Pierre Elliott Trudeau International Airport,45.467837,-73.742294,118,NA,CA,CA-QC,Montréal,yes,CYUL,YUL,CYUL,YUL,http://www.admtl.com/passager/Home.aspx,https://en.wikipedia.org/wiki/Montréal–Trudeau_International_Airport,"YMQ, Dorval Airport"
\end{lstlisting}

Para cada aeropuerto con código IATA se procede a crear un conjunto de reglas para que el \textit{pipeline} pueda
reconocer con éxito la entidad \texttt{AIRPORT\_CODE} mencionando todos los códigos de aeropuerto existentes.
Se creó un archivo \texttt{airport.rules} que asocia tanto el código IATA como la combinación de código IATA seguido del
nombre de la ciudad donde se ubica el aeropuerto.
Retomando el ejemplo de Montreal, el contenido del archivo de reglas es el siguiente:

\begin{lstlisting}[
    style=basecodestyle,
    caption={Reglas de \texttt{airport.rules} para el aeropuerto YUL.},
    label={lst:airport-rules-yul}
]
YUL	AIRPORT_CODE	O,MISC,ORGANIZATION	2.0
YUL Montréal	AIRPORT_CODE	O,MISC,ORGANIZATION	2.0
\end{lstlisting}

La estructura de los archivos de reglas consiste en hasta 4 columnas separadas por un \textit{tab}.
En la primer columna se indica la cadena a corresponder, luego se indica  el tipo de categoría a asociar.
En la tercer columna se indican las categorías reconocidas que pueden ser reemplazadas (en el ejemplo estamos indicando
que si fue reconocido como \texttt{MISC} o como \texttt{ORGANIZATION} entonces es más importante \texttt{AIRPORT\_CODE}
y por lo tanto debería reconocerlo así.
Por último con el valor \texttt{2.0} estamos indicando la prioridad, que resultá útil para el desempate cuando dos
reglas aplican para una misma cadena.
De esta forma CoreNLP entiende que frente al tag ``YUL'' lo debe reconocer como un nuevo tipo \texttt{AIRPORT\_CODE}.

Debido a que CoreNLP puede reconocer a una cadena ``YUL Montréal'' como una única entidad, y anotarla como
\texttt{ORGANIZATION}, entonces necesitamos de una segunda regla para cada aeropuerto donde se ubica primero el código
IATA y luego el nombre de la ciudad correspondiente.

La particularidad de los nombres de los aeropuertos es que el modelo puede fácilmente con el nombre de una persona,
por ejemplo del pasajero del vuelo.
Siguiendo con el mismo ejemplo, para la cadena ``Charles de Gaulle Airport'' el modelo lo anota como \texttt{PERSON}.
Es por eso que para cada aeropuerto se crean hasta tres reglas adicionales indicadas en el archivo
\texttt{airport-names.rules} donde para un nombre proveniente del dataset como ``Montreal / Pierre Elliott Trudeau
International Airport'' se incluye no sólo la cadena exacta, sino también la misma sin la palabra ``Airport'' y otra sin
la cadena ``International Airport''.
De esta forma, para el aeropuerto en cuestión, se obtienen las siguientes reglas:

\begin{lstlisting}[
    style=basecodestyle,
    caption={Reglas de \texttt{airport-names.rules} para el aeropuerto YUL.},
    label={lst:airport-names-rules-yul}
]
Montreal / Pierre Elliott Trudeau	AIRPORT_NAME	O,MISC,ORGANIZATION,PERSON	2.0
Montreal / Pierre Elliott Trudeau International	AIRPORT_NAME	O,MISC,ORGANIZATION,PERSON	2.0
Montreal / Pierre Elliott Trudeau International Airport	AIRPORT_NAME	O,MISC,ORGANIZATION,PERSON	2.0
\end{lstlisting}

A diferencia del fragmento de reglas anterior aquí es necesario precisar que el tipo de entidad \texttt{PERSON} es
reemplazable por una nueva entidad \texttt{AIRPORT\_NAME}.
En principio para el procesamiento de los correos de vuelos no es necesario extraer el nombre del aeropuerto ya que los
identificadores IATA anotados como \texttt{AIRPORT\_CODE} son únicos.
De todas formas la identificación es útil para que el nombre del aeropuerto no sea identificado como \texttt{PERSON} y
así se facilita la tarea de identificar el nombre del pasajero de la reserva.

Por último se creó una regla de \texttt{TokensRegex} (\url{https://nlp.stanford.edu/software/tokensregex.html}) donde se
rescata al aeropuerto cuando el texto visible queda con un formato como ``EZE Buenos Aires Departing at'' o ``EZE Buenos
Aires Arriving at:''.
Para ello, en el archivo \texttt{contextual-airport.tokensregex.rules} se definió lo siguiente:

\begin{lstlisting}[
    style=jsonlikecodestyle,
    caption={Regla de \texttt{contextual-airport.tokensregex.rules}.},
    label={lst:contextual-airport-tokensregex}
]
ner = { type: "CLASS", value: "edu.stanford.nlp.ling.CoreAnnotations$NamedEntityTagAnnotation" }
{
  ruleType: "tokens",
  pattern: ((([{word:/[A-Z]{3}/}] [{word:/[A-Z][A-Za-z'.-]*/}] [{word:/[A-Z][A-Za-z'.-]*/}]?)) [{word:/Departing|Arriving/}] [{word:/at/}] [{word:/:/}]?),
  action: Annotate($2, ner, "AIRPORT_CODE"),
  result: "AIRPORT_CODE"
}
\end{lstlisting}

Como se puede apreciar se define un patrón donde primero aparece el código IATA de tres letras, seguido de 0 o más
palabras y luego las palabras exactas ``Departing'' o ``Arriving''.
Esta regla permite identificar al aeropuerto según el contexto en el que se ubica, pudiendo anotar únicamente el código
IATA como \texttt{AIRPORT\_CODE} y descartando el resto (en la parte de \texttt{action} se indica
\texttt{Annotate(\$2,...)} para anotar el segundo grupo capturado, esto es, ``EZE Buenos Aires''.

A pesar de contar con reglas para la totalidad de códigos de aeropuertos existentes esto trae aparejado el problema de
los falsos positivos.
Una sigla en inglés como ``MAY'' que corresponde a un \textit{modal verb} será identificado incorrectamente como el
``Clarence A. Bain Airport'' en las Bahamas.
Del mismo modo la cadena ``MAX'' utilizada para indicar el máximo peso aceptado por la aerolínea en el equipaje ahora
será identificado como el ``Ouro Sogui Airport'' ubicado en Senegal.
Para estos casos extremos identificados, que siempre corresponden a aeropuertos pequeños, se decidió utilizar la
propiedad \texttt{ner.additional.regexner.commonWords} para indicarle a CoreNLP las palabras que deben ser tomadas como
palabras comunes del idioma inglés y así excluirlas de cualquier identificación.
Las palabras definidas en la implementación son las siguientes:

\begin{lstlisting}[
    style=basecodestyle,
    caption={Palabras comunes excluidas por \texttt{regexner}.},
    label={lst:regexner-common-words}
]
GET
MAY
VOL
API
GST
HST
MAX
\end{lstlisting}

\subsubparagraph{Códigos de Vuelos y Aerolíneas}

Anteriormente se comentó el preprocesamiento que se realiza antes de iniciar el pipeline de CoreNLP y que involucra a
los códigos de vuelo definidos por IATA.
Ahora es momento de obtener el token ``AC8633'' y reconocerlo como una entidad de tipo \texttt{FLIGHT\_NUMBER}.
Para ello se utilizó el \textit{dataset} público de aerolíneas de OpenFlights (\url{https://openflights.org/data.php})
que ofrece un CSV con más de 6.000 registros que representan a todas las aerolíneas del mundo.
Por ejemplo:

\begin{lstlisting}[
    style=basecodestyle,
    caption={Registro de ejemplo del \textit{dataset} OpenFlights.},
    label={lst:openflights-example}
]
C1,C2,C3,C4,C5,C6,C7,C8
330,Air Canada,\N,AC,ACA,AIR CANADA,Canada,Y
\end{lstlisting}

Para cada aerolínea se cuenta con el código de dos letras de IATA (``AC''), el nombre y el país de origen de la misma,
entre otros.
Con esta información se puede crear un nuevo conjunto de reglas que modele el comportamiento de encontrar todos los
códigos de aerolínea (compuesto por dos caracteres alfanuméricos) seguidos de entre 2 y 4 dígitos que corresponden al
número de vuelo.
Para la aerolínea en cuestión el archivo de reglas es el siguiente:

\begin{lstlisting}[
    style=basecodestyle,
    caption={Regla de \texttt{airlines.rules} para Air Canada.},
    label={lst:airlines-rules-air-canada}
]
AC\d{2,4}	FLIGHT_NUMBER	O,MISC	2.0
\end{lstlisting}

donde el nuevo tipo de entidad \texttt{FLIGHT\_NUMBER} tiene prioridad por sobre una entidad de tipo \texttt{MISC}.
Este conjunto se reglas se ubica en el archivo \texttt{airlines.rules}.

Utilizando el mismo \textit{dataset} se crea otro conjunto de reglas, ubicado en el archivo \texttt{airline-names.rules}
para definir el nuevo tipo de entidad \texttt{AIRLINE\_NAME} y que tenga prioridad por sobre tipos \texttt{MISC} y
\texttt{ORGANIZATION}.

\begin{lstlisting}[
    style=basecodestyle,
    caption={Regla de \texttt{airline-names.rules} para Air Canada.},
    label={lst:airline-names-rules-air-canada}
]
AIR CANADA	AIRLINE_NAME	O,MISC,ORGANIZATION	2.0
\end{lstlisting}

Nuevamente, si bien no es necesario extraer el nombre de la aerolínea (ya que la primera porción del código de vuelo
es un identificador único), el hecho de anotar un nombre como nombre de aerolínea ayuda a que existen menos partes del
contenido sin anotar y que podrían ser candidatas a identificaciones incorrectas.

\subsubparagraph{Código de Reserva (PNR)}

Otro campo relevante para la extracción es el código de reserva.
Aquí también IATA define el \texttt{PNR} o \textit{Passanger Name Record}
(\url{https://en.wikipedia.org/wiki/Passenger_name_record}) y consiste en 6 caracteres alfanuméricos.
Para dificultad de la extracción se recuerda que el código de vuelo (como ``AC8633'') ya está cumpliendo con el formato
de \texttt{PNR}.
Es por eso que aquí la estrategia para la extracción es distinta y consiste en identificar el \texttt{PNR} por contexto.
Se establecen un conjunto de reglas \texttt{TokensRegex} para extraer el \texttt{PNR} cuando aparece luego de frases como
``Confirmation Number'', ``Booking Code'', ``Booking Reference'', entre otras.
En el archivo \texttt{pnr.tokensregex.rules} se indican seis reglas.
El principio del archivo es el siguiente:

\begin{lstlisting}[
    style=jsonlikecodestyle,
    caption={Reglas iniciales de \texttt{pnr.tokensregex.rules}.},
    label={lst:pnr-tokensregex-rules}
]
ner = { type: "CLASS", value: "edu.stanford.nlp.ling.CoreAnnotations$NamedEntityTagAnnotation" }
$PNR_CODE = "/[A-Z0-9]{6}/"
$SEPARATOR = "/[:-]/"
{
  ruleType: "tokens",
  pattern: ([{word:/[Bb]ooking/}] [{word:/[Rr]eference/}] [{word:$SEPARATOR}]? ([{word:$PNR_CODE}])),
  action: Annotate($1, ner, "PNR_CODE"),
  result: "PNR_CODE"
}
\end{lstlisting}

Con estas reglas se define un nuevo tipo de entidad \texttt{PNR\_CODE} listo para su extracción.

\subsubparagraph{SUTime: Fechas y horas}

En un mismo mail de confirmación de reserva pueden aparecer múltiples fechas y horas correspondientes por ejemplo al
despegue o aterrizaje de cualquiera de los tramos de la reserva.
Además es común que se indique la fecha en que se realizó la confirmación o incluso la fecha límite hasta donde se
pueden realizar modificaciones a la reserva sin costo adicional.
Es por eso que en esta etapa de anotación sólo resulta útil anotar todos los campos correspondientes a fecha y hora como
\texttt{TIME} y la extracción con contexto se realizará en la siguiente etapa.
Solo así podremos asociar por ejemplo la fecha y hora de despegue al primer tramo de una reserva y no al segundo.

Las reglas por defecto fueron suficientes para la mayoría de los correos electrónicos, ya que CoreNLP provee la
biblioteca SUTime (\url{https://nlp.stanford.edu/software/sutime.html}) que permite reconocer por ejemplo fechas con
múltiples formatos como ``31-12-2026'', ``31/12/2026'' o ``2026-31-12''.

Sólo fue necesaria una regla \texttt{TokensRegex} extra:

\begin{lstlisting}[
    style=jsonlikecodestyle,
    caption={Regla adicional para fechas en boarding passes.},
    label={lst:sutime-boarding-pass-rule}
]
ner = { type: "CLASS", value: "edu.stanford.nlp.ling.CoreAnnotations$NamedEntityTagAnnotation" }
{
  ruleType: "tokens",
  pattern: (([{word:/\d{2}[A-Z]{3}\d{2}/}])),
  action: Annotate($1, ner, "DATE"),
  result: "DATE"
}
\end{lstlisting}

para que en mails correspondientes a \textit{boarding passes}, una cadena como ``22SEP25'' sea interpretada como
``22/09/2025''.

\subsubparagraph{Marcadores de tramos}

Por último se define un archivo \texttt{segment-context.rules} con las siguientes dos reglas:
\begin{lstlisting}[
    style=basecodestyle,
    caption={Reglas de \texttt{segment-context.rules}.},
    label={lst:segment-context-rules}
]
Departing	FLIGHT_SEGMENT_MARKER	O,MISC,ORGANIZATION	2.0
Arriving	FLIGHT_SEGMENT_MARKER	O,MISC,ORGANIZATION	2.0
\end{lstlisting}

Si bien en la siguiente etapa se podría buscar la palabra exacta ``Departing'' es mejor anotarla como una nueva entidad
de tipo \texttt{FLIGHT\_SEGMENT\_MARKER} para buscar múltiples entidades de un tipo en particular y así identificar uno
o más tramos de vuelos dentro de una misma reserva.

\paragraph{Extracción de información}

\subparagraph{Extracción de campos únicos: PNR y Pasajero}

\subparagraph{Extracción de campos múltiples: Segmentos de vuelo}

\subsubparagraph{Cuatro Layouts posibles}

\subsection{Resultados Experimentales}

\section{Conclusiones y Trabajo a Futuro}
